\documentclass[11pt]{article}

\usepackage[margin=1in]{geometry}
\usepackage[protrusion=true,expansion=false]{microtype}
\usepackage{booktabs}
\usepackage{graphicx}
\usepackage{amsmath}
\usepackage{array}
\usepackage{longtable}
\usepackage{threeparttable}
\usepackage{caption}
\usepackage{float}
\usepackage{pdflscape}
\usepackage{xcolor}
\usepackage{xspace}
\usepackage[hidelinks]{hyperref}

\graphicspath{{outputs/final/figures/}}
% Generated by src/analysis/run_pipeline.py; do not edit.
\newcommand{\RawRows}{3,742\xspace}
\newcommand{\PreferredRows}{3,567\xspace}
\newcommand{\Targets}{1,834\xspace}
\newcommand{\ParentEntities}{1,830\xspace}
\newcommand{\GlobalEntities}{1,818\xspace}
\newcommand{\UINs}{186\xspace}
\newcommand{\ObservedRoots}{184\xspace}
\newcommand{\Parents}{28\xspace}
\newcommand{\Jurisdictions}{122\xspace}
\newcommand{\NonrootEdges}{1,650\xspace}
\newcommand{\ExpectedLinks}{1,500\xspace}
\newcommand{\OtherLevelLinks}{72\xspace}
\newcommand{\CompletePaths}{1,687\xspace}
\newcommand{\CompletePathPct}{92.0\xspace}
\newcommand{\MissingParentEdges}{78\xspace}
\newcommand{\MissingParentNodes}{18\xspace}
\newcommand{\CrossBorderEdges}{951\xspace}
\newcommand{\CrossBorderPct}{57.6\xspace}
\newcommand{\CrossBorderEqualPct}{58.7\xspace}
\newcommand{\CrossBorderLooLow}{56.1\xspace}
\newcommand{\CrossBorderLooHigh}{61.8\xspace}
\newcommand{\ModalUinPct}{59.7\xspace}
\newcommand{\ModalUinCount}{1,095\xspace}
\newcommand{\ModalUinEqualPct}{62.1\xspace}
\newcommand{\ModalUinMedianPct}{61.2\xspace}
\newcommand{\ObservedSubtreeMedian}{43.2\xspace}
\newcommand{\NamedNodeSubtreeMedian}{50.8\xspace}
\newcommand{\DagDominatorMedian}{41.7\xspace}
\newcommand{\ConcentrationParents}{24\xspace}
\newcommand{\ObservedIntermediaries}{403\xspace}
\newcommand{\DutchGatewayMean}{34.8\xspace}
\newcommand{\DutchGatewayMedian}{12.0\xspace}
\newcommand{\DutchGatewayEqualMean}{37.1\xspace}
\newcommand{\DutchGatewayLooLow}{21.2\xspace}
\newcommand{\DutchGatewayLooHigh}{39.7\xspace}
\newcommand{\DutchGatewayRobustMean}{37.8\xspace}
\newcommand{\DutchUpstreamPct}{29.4\xspace}
\newcommand{\DutchUpstreamEqualPct}{20.9\xspace}
\newcommand{\TopThreeCount}{1,140\xspace}
\newcommand{\TopThreePct}{62.2\xspace}
\newcommand{\TopThreeEqualPct}{46.7\xspace}
\newcommand{\TopThreeNoFundsPct}{59.8\xspace}
\newcommand{\DeclaredCentrePct}{61.8\xspace}
\newcommand{\DeclaredCentreEqualPct}{57.2\xspace}
\newcommand{\DeclaredCentreNodePct}{38.0\xspace}
\newcommand{\SourceTargets}{847\xspace}
\newcommand{\ParsedTargets}{444\xspace}
\newcommand{\ReadyRows}{560\xspace}
\newcommand{\SignPlausibleRows}{521\xspace}
\newcommand{\PlValidRows}{105\xspace}
\newcommand{\ReadyTargets}{265\xspace}
\newcommand{\ParentsNoReady}{5\xspace}
\newcommand{\RepeatedClusters}{51\xspace}
\newcommand{\RepeatedRows}{104\xspace}
\newcommand{\ReadyRepeatedRows}{70\xspace}
\newcommand{\EvidenceReuseClusters}{33\xspace}
\newcommand{\BroadUnflaggedRows}{460\xspace}
\newcommand{\MissingUnitsReady}{516\xspace}
\newcommand{\DepthTwoPlusCount}{871\xspace}
\newcommand{\DepthTwoPlusPct}{47.5\xspace}
\newcommand{\DepthFivePlusCount}{342\xspace}
\newcommand{\DepthFivePlusPct}{18.6\xspace}
\newcommand{\DepthFiveEqualPct}{7.3\xspace}
\newcommand{\GraphDepthFivePct}{7.4\xspace}
\newcommand{\DepthMismatchCount}{236\xspace}
\newcommand{\DepthMismatchPct}{14.0\xspace}
\newcommand{\MothersonMismatchShare}{88.1\xspace}
\newcommand{\DeepManufacturing}{242\xspace}
\newcommand{\ManufacturingTargets}{508\xspace}
\newcommand{\DeepManufacturingTopTwo}{240\xspace}
\newcommand{\DeepManufacturingTopTwoPct}{99.2\xspace}
\newcommand{\ZeroStakeCount}{406\xspace}
\newcommand{\ZeroStakePct}{24.6\xspace}
\newcommand{\OwnershipPooledGap}{26.4\xspace}
\newcommand{\OwnershipEqualGap}{2.3\xspace}
\newcommand{\OwnershipNoRelianceGap}{6.1\xspace}
\newcommand{\OwnershipLevelOne}{56.4\xspace}
\newcommand{\OwnershipLevelTwoPlus}{82.8\xspace}
\newcommand{\FinancialRawDepthGap}{-6.6\xspace}
\newcommand{\FinancialFeDepthGap}{-2.4\xspace}
\newcommand{\FinancialLooLow}{-6.6\xspace}
\newcommand{\FinancialLooHigh}{-0.8\xspace}
\newcommand{\HoldingNameCount}{177\xspace}
\newcommand{\HoldingNameHolderPct}{64.4\xspace}
\newcommand{\OtherNameHolderPct}{17.6\xspace}
\newcommand{\HoldingNameFeGap}{36.3\xspace}
\newcommand{\TataCommRawRows}{550\xspace}
\newcommand{\TataCommEntities}{49\xspace}
\newcommand{\MothersonRawRows}{332\xspace}
\newcommand{\MothersonEntities}{309\xspace}


\definecolor{paperblue}{HTML}{356A9A}
\hypersetup{
  pdftitle={Mapping the First Hop},
  pdfauthor={Canonical research synthesis},
  colorlinks=true,
  linkcolor=paperblue,
  urlcolor=paperblue
}
\captionsetup{font=small,labelfont=bf}
\setcounter{secnumdepth}{3}
\setcounter{tocdepth}{2}

\newcommand{\pct}{\%}
\xspaceaddexceptions{\pct}
\newcommand{\pp}{percentage points}
\newcommand{\sample}{the supplied \Parents-parent dataset}
\newcommand{\FloatBarrier}{\clearpage}

\title{\textbf{Mapping the First Hop: Denominators, Hierarchy, and Gateway
Structure in Reported Indian Foreign Ownership Networks}}
\author{Canonical research synthesis}
\date{August 2026}

\begin{document}

\maketitle

\begin{abstract}
This paper reconstructs the reported foreign ownership networks of \Parents
Indian parent buckets in a supplied research dataset. The source is not a firm
cross-section: \RawRows\ candidate source rows reduce to \PreferredRows
preferred target-years, \Targets\ structural target occurrences,
\ParentEntities\ parent-scoped normalized entity candidates, and
\GlobalEntities\ global normalized candidates. Keeping these denominators
separate changes both rankings and interpretation. The modal registration or
mapping channel contains \ModalUinPct\pct\ of targets in the pooled data and
\ModalUinEqualPct\pct\ under equal-parent weighting. Among
\ConcentrationParents\ sufficiently large parent groups, the median largest
observed subtree contains \ObservedSubtreeMedian\pct\ of normalized entities;
the corresponding strict directed-acyclic-graph dominator estimate is
\DagDominatorMedian\pct. First-hop gateways in the Netherlands have a mean of
\DutchGatewayMean\ downstream entities, compared with a parent-equal mean of
\DutchGatewayEqualMean, and a leave-one-parent-out range of
\DutchGatewayLooLow--\DutchGatewayLooHigh. Direct relations are also commonly
international: \CrossBorderEdges/\NonrootEdges\ nonroot edges
(\CrossBorderPct\pct) cross a jurisdiction boundary, with an equal-parent
estimate of \CrossBorderEqualPct\pct. These findings survive entity
de-duplication, parent balancing, leave-one-parent-out checks, and alternative
network definitions. Other apparent patterns do not. Reported level and
reconstructed graph distance disagree for \DepthMismatchCount/\CompletePaths
complete paths, and pooled depth, ownership, sector, and financial patterns are
strongly affected by parent composition or measurement limitations. The paper
therefore establishes a descriptive topology for this dataset, not causal,
legal, tax, regulatory, or misconduct conclusions.
\end{abstract}

\noindent\textbf{Keywords:} ownership networks; corporate hierarchy;
subsidiary geography; entity resolution; robustness; denominator audit

\vspace{0.8em}
\noindent\textbf{Scope note.} This first canonical version deliberately uses
only evidence validated inside the supplied data and code. It does not add an
external literature review or institutional interpretation, because those
sources were not part of the adjudicated evidence base. Such context is a task
for a separately sourced revision.

\tableofcontents
\newpage

\section{Introduction}

Datasets on international corporate structure often combine objects that look
deceptively interchangeable. These include source records, reporting years,
registered investments, named subsidiaries, parent relations, and complete
ownership paths. Treating any one of these as ``the number of firms''
can mechanically change which parent groups or jurisdictions dominate a
result. Hierarchy adds a second problem. A location can be common because many
entities reside there, because a smaller number of upstream entities have many
descendants, or because one unusually large parent group repeats the same
route. Those are different empirical statements.

This paper asks a deliberately bounded question: how do the parent groups in
\sample\ organize their reported foreign ownership networks across mapping
channels, hierarchy, and jurisdictions, and which patterns remain after
correcting units, duplicates, parent weights, missing nodes, and graph
definitions? The goal is measurement-corrected description. The data do not
identify why a structure exists, whether it is current, whether a named entity
has a particular legal function, or whether any route has a tax or regulatory
purpose.

The analysis makes five contributions within that boundary. First, it builds
an explicit denominator ladder from source observations to target, entity,
edge, path, mapping-channel, and parent units. Second, it reconstructs a
parent-scoped directed ownership graph while retaining named-but-unscraped
parents and their supplied countries. Third, it separates resident geography
from intermediary and descendant exposure. Fourth, it reports pooled,
equal-parent, leave-one-parent-out, and alternative-definition estimates side
by side. Fifth, it treats failed hypotheses as results: several initially
striking depth, ownership, sector, and financial relationships do not support a
general interpretation after the relevant correction.

The central evidence is as follows. Parent networks are heterogeneous along
dimensions that cannot be collapsed into a single defensible complexity
ranking. Many are concentrated behind one mapping channel or one large
observed subtree. Dutch and Mauritian first-hop entities have greater
downstream amplification than their resident counts alone would suggest, while
the United States is primarily prominent as a resident location. A majority
of reported nonroot edges cross jurisdictions, and that result is unusually
stable to parent balance. In contrast, reported deep hierarchy is concentrated
in a few parent groups and differs materially from reconstructable graph
distance. The financial layer is sparse and extraction-contaminated, so it
supports an audit of measurement quality rather than comparative financial
performance.

The remainder of the paper defines the units and graph, documents the
denominator correction, presents parent and jurisdictional topology, examines
depth and financial limitations, and concludes with an explicit evidence
boundary.

\section{Data and analytical units}
\label{sec:data}

\subsection{Source structure and scope}

The immutable input is a Stata file containing candidate financial-statement
sources attached to reported foreign targets. The file has \RawRows\ rows and
\Jurisdictions\ standardized jurisdiction labels across \Parents\ supplied
ultimate-parent buckets. Rows are not unique companies. A target can recur for
multiple fiscal years and can have multiple candidate sources within a year.
The upstream preferred-source flag selects \PreferredRows\ target-year rows,
but structural fields are invariant within each target identifier. The
structural cross-section therefore contains \Targets\ target occurrences.

The sample is not accompanied by a sampling frame for all Indian outward
investment. No result below is generalized beyond the supplied parent buckets.
Similarly, the file provides reported names and ownership relationships, not a
complete legal-entity registry or verified historical ownership charts.

\begin{table}[tbp]
\centering
\caption{The canonical denominator ladder}
\label{tab:denominators}
\small
\begin{tabular}{lrl}
\toprule
Analytical unit & Count & Definition \\
\midrule
Raw source rows & 3,742 & source candidate \\
Preferred target-year rows & 3,567 & target\_id $\times$ fiscal\_year \\
Structural target occurrences & 1,834 & target\_id \\
Parent-scoped normalized entity candidates & 1,830 & parent $\times$ normalized name-country \\
Global normalized entity candidates & 1,818 & normalized name-country \\
UINs & 186 & uin \\
Observed level-0 targets & 184 & target occurrence \\
Ultimate-parent buckets & 28 & source parent label \\
Complete observed paths & 1,687 & target path \\
\bottomrule
\end{tabular}

\begin{minipage}{0.94\textwidth}
\footnotesize\vspace{0.4em}
\textit{Notes:} Counts are generated from the immutable source by the
canonical pipeline. A target occurrence is one reported ownership-path target.
Normalized entities are conservative name--country candidates, not verified
legal identities. A complete observed path reaches the Indian parent through
observed entity nodes.
\end{minipage}
\end{table}

\begin{figure}[tbp]
\centering
\includegraphics[width=\textwidth]{figure_01_denominators_and_depth.pdf}
\caption{Analytical denominators and two measures of hierarchy depth}
\label{fig:denominators}
\begin{minipage}{0.94\textwidth}
\footnotesize\textit{Notes:} Panel (a) uses a logarithmic count scale because
the units differ sharply in size. Panel (b) compares the supplied reported
level for all targets with graph distance among complete observed paths; these
are not interchangeable measures.
\end{minipage}
\end{figure}

\subsection{Targets, normalized entities, and mapping channels}

The primary path and exposure unit is the target identifier. For within-parent
architecture, entity names are normalized conservatively: Unicode is reduced
to ASCII, case and punctuation are removed, ampersands are standardized, and
the supplied country is spelling-harmonized. A parent-scoped entity key is the
combination of ultimate-parent bucket, normalized name, and standardized
country. This yields \ParentEntities\ candidates, only slightly fewer than the
\Targets\ targets. Dropping the parent component yields \GlobalEntities
global candidates. The difference is small in aggregate, but the estimands are
conceptually distinct: a shared venture may be one global entity while still
representing exposure in more than one parent network.

The file also contains \UINs\ UIN values. They are retained as supplied and
treated only as registration or mapping channels. The analysis does not infer
the semantics of their substrings, a registration date, or legal-entity
identity. \ObservedRoots\ UINs have an observed level-0 target, so two mapping
channels lack an observed first-hop entity in this file.

\subsection{Financial observations}

Financial variables remain at the preferred target-year level. Coverage gates
are reported separately from network construction. A target has a located
source if any source row carries the supplied source-found indicator, a parsed
balance sheet if equity is numeric in any row, and a valuation-ready record if
the upstream ready flag is present at least once. Absolute monetary values are
not pooled across currencies or unspecified units.

\section{Ownership-network construction and estimands}
\label{sec:methods}

\subsection{Edge resolution}

Each structural target supplies an entity name and country, reported level,
immediate-parent name and country, UIN, and ultimate-parent bucket. Candidate
parents are formed within the same ultimate-parent bucket from conservative
name aliases. Candidate matches are ranked lexicographically by whether they
occupy the expected preceding level, whether the supplied parent country
matches, whether the original names match exactly, and whether the UIN matches.
The canonical build accepts no fuzzy match.

The result is one occurrence-level edge for each target. Level-0 targets link
to an explicit Indian ultimate-parent node. Below level 0, \ExpectedLinks\ of
\NonrootEdges\ relations link to an observed entity at the expected prior
level, \OtherLevelLinks\ link to an observed entity at another reported level,
and \MissingParentEdges\ terminate at one of \MissingParentNodes\ normalized
named-but-unscraped parent nodes. These unresolved nodes retain their supplied
country. Discarding the country would understate upstream exposure for some
groups; treating the node as observed would overstate path completeness.

Occurrence edges are collapsed to a logical entity graph for subtree and
gateway calculations. The parent-specific normalized graphs are directed and
acyclic under the accepted links. The graph preserves alternative reported
parents where they exist rather than forcing a tree.

\subsection{Paths and depth}

For every target, recursion follows observed parents toward the ultimate
parent. A path is complete only when it reaches that root through observed
entities. This succeeds for \CompletePaths/\Targets\ targets
(\CompletePathPct\pct). A path that encounters a named missing node remains
useful for country exposure up to that node but is not declared complete.

The supplied reported level and reconstructed graph distance are stored
separately. Reported level may encode a hierarchy with nodes missing from the
target table. Conversely, cross-channel name matching can create a shorter
observed route than the reported sequence. Without external ownership charts,
neither measure is relabeled as ``true'' depth.

\subsection{Network concentration}

Four concentration quantities answer different questions.

\begin{itemize}
  \item The \emph{modal-UIN share} is the fraction of a parent's targets
  assigned to its largest UIN channel.
  \item The \emph{largest observed subtree} counts normalized entity
  descendants below the observed entity with the largest descendant set,
  divided by the other entities in that parent network.
  \item The \emph{named-node-inclusive subtree} also allows an unresolved
  named parent to be the top node, while still not calling the path complete.
  \item The \emph{strict DAG dominator} removes an observed node and counts the
  entity nodes no longer reachable from the ultimate parent. It therefore
  respects alternative routes that bypass the candidate node.
\end{itemize}

The subtree summaries use the \ConcentrationParents\ parent groups with at
least fifteen targets. That threshold prevents very small groups from
determining a median while retaining most of the structural sample. Subtree
language describes the mapped topology. It does not establish legal control,
cash-flow dependence, or an operational chokepoint.

\subsection{Geographic roles}

The analysis distinguishes three jurisdiction roles. \emph{Residence} counts
normalized entities located in a jurisdiction. An \emph{observed
intermediary} is a unique observed entity that appears as the parent of at
least one logical edge. A jurisdiction is \emph{strictly upstream} of a target
if it appears above the terminal target in the reconstructed path; the target's
own residence is excluded from its exposure set. India is excluded as the
ultimate root in foreign-jurisdiction role and cover calculations.

A \emph{gateway} is an observed level-0 entity. Gateway amplification is the
number of normalized entity descendants below it. Jurisdiction means weight
gateways; parent-equal means first average gateways within parent and then
average parents. A robustness subset retains gateways for which no single
direct-child branch contains a majority of descendants.

An edge is cross-border when the standardized child and immediate-parent
country labels differ. The cross-border statistic excludes level-0 links to
India, which would otherwise be cross-border by construction.

\subsection{Weighting and robustness}

For a binary target outcome $y_i$ and parent $p(i)$, the pooled estimand is
$N^{-1}\sum_i y_i$. The equal-parent estimand is

\[
  \frac{1}{P}\sum_{p=1}^{P}\left(\frac{1}{N_p}\sum_{i:p(i)=p} y_i\right).
\]

The former describes target-weighted exposure in this file; the latter gives
each parent bucket the same weight. Leave-one-parent-out estimates repeat the
relevant pooled or fixed-effect calculation after excluding each parent in
turn. Where units differ---entities, edges, paths, gateways, or parents---the
paper labels the denominator explicitly rather than comparing percentages as
though they shared an estimand.

\section{Denominators and parent architecture}
\label{sec:architecture}

\subsection{Raw-row weighting changes the sample composition}

The denominator correction is empirically first order. Tata Communications
contributes \TataCommRawRows\ source rows but only \TataCommEntities
parent-scoped normalized entities. Motherson contributes \MothersonRawRows
rows but \MothersonEntities\ normalized entities. Thus the former is largest
by source rows while the latter is largest by network entities. Any statistic
computed directly on source rows would overweight targets with longer or more
duplicated financial-source histories.

The closeness of the target and normalized-entity totals is reassuring but
does not make the units interchangeable. Four parent-scoped alias collapses
are enough to matter for exact graph topology, and cross-parent shared names
have a different interpretation from within-parent duplicates. The canonical
tables therefore retain target and entity counts side by side.

\subsection{Architecture is multidimensional}

Figure~\ref{fig:parent-architecture} shows that group size, geographic breadth,
reported depth, and channel concentration do not move together. Table
\ref{tab:parent-examples} gives illustrative contrasts. Motherson is the
largest entity network in the sample and has the greatest reported maximum
depth, while Wipro spans more jurisdiction labels. Reliance Industries is
large but shallow in reported levels. Tata Communications is much smaller as
an entity network than its source-row count suggests. ONGC Videsh has many
UIN channels relative to its target count and a shallow observed graph.

These comparisons do not support a one-dimensional complexity index. A group
can be broad but shallow, deep but concentrated behind one entry channel, or
large because one subtree contains many acquired or portfolio entities.
Calling any one of those structures ``more complex'' would impose an
unvalidated weighting across substantively different axes.

\begin{figure}[tbp]
\centering
\includegraphics[width=0.91\textwidth]{figure_02_parent_architecture.pdf}
\caption{Parent groups combine size, geographic breadth, and depth differently}
\label{fig:parent-architecture}
\begin{minipage}{0.9\textwidth}
\footnotesize\textit{Notes:} Each marker is one supplied ultimate-parent
bucket. Marker color is maximum reported level; marker area increases with the
within-parent modal-UIN target share. Entity count uses the parent-scoped
normalized name--country key.
\end{minipage}
\end{figure}

\begin{table}[tbp]
\centering
\caption{Illustrative parent architectures}
\label{tab:parent-examples}
\scriptsize
\resizebox{\textwidth}{!}{\begin{tabular}{lrrrrrrr}
\toprule
Parent bucket & Raw rows & Entities & Countries & Reported max & Graph max & Modal UIN (\%) & Complete paths (\%) \\
\midrule
Motherson & 332 & 309 & 45 & 12 & 8.0 & 44.7 & 100 \\
Wipro & 504 & 212 & 51 & 8 & 7.0 & 50.2 & 88.7 \\
Reliance Industries & 277 & 196 & 22 & 2 & 2.0 & 34.7 & 100 \\
Hindalco & 81 & 81 & 25 & 12 & 6.0 & 95.1 & 46.9 \\
Tata Communications & 550 & 49 & 32 & 2 & 2.0 & 68.0 & 100 \\
ONGC Videsh & 41 & 41 & 18 & 2 & 2.0 & 29.3 & 100 \\
\bottomrule
\end{tabular}
}
\begin{minipage}{0.98\textwidth}
\footnotesize\vspace{0.4em}
\textit{Notes:} Graph maximum is reconstructed distance among complete paths.
Modal UIN is the largest UIN's share of target occurrences. Complete paths are
reported as a percentage of the parent's targets. Rows are selected to
illustrate contrasting architectures, not to rank firms.
\end{minipage}
\end{table}

\FloatBarrier

\section{Dominant channels and large subtrees}
\label{sec:channels}

The largest UIN within each parent accounts for \ModalUinCount/\Targets
targets, or \ModalUinPct\pct\ pooled. Giving every parent equal weight produces
\ModalUinEqualPct\pct; the median parent is \ModalUinMedianPct\pct, and the
pooled leave-one-parent-out interval is narrow (Figure~\ref{fig:robustness}).
Channel concentration is therefore not only an artifact of the largest parent
bucket. At the same time, UIN is an administrative or mapping unit in this
analysis, not a verified company or control node.

Graph concentration yields a complementary result. Across the
\ConcentrationParents\ larger groups, the median largest observed subtree
contains \ObservedSubtreeMedian\pct\ of the other normalized entities. The
median changes only modestly to \DagDominatorMedian\pct\ under the strict DAG
dominator definition, which allows alternative paths to bypass the removed
node. When named-but-unscraped parent nodes may sit at the top of a subtree,
the median rises to \NamedNodeSubtreeMedian\pct.

Figure~\ref{fig:channels} makes the conceptual separation visible. UIN and
subtree shares are often close, but not always. A mapping channel can contain
many targets whose observed graph is split, while a linked subtree can span
more than one UIN. Missing parents can also reverse the apparent concentration
of a particular group. The defensible result is that dominant mapped channels
and large observed subtrees are common in this sample. It is not that every
network has one legally dispositive company.

\begin{figure}[p]
\centering
\includegraphics[width=\textwidth]{figure_03_channel_concentration.pdf}
\caption{Channel concentration under alternative graph definitions}
\label{fig:channels}
\begin{minipage}{0.96\textwidth}
\footnotesize\textit{Notes:} The sample is parent groups with at least fifteen
targets. Subtree shares use normalized entity candidates and exclude the top
node from the denominator. The strict DAG measure counts entities detached
from the ultimate root after removing the candidate node. The hollow marker is
the target-weighted modal-UIN share and therefore uses a different unit.
\end{minipage}
\end{figure}

\FloatBarrier

\section{Geography and intermediary structure}
\label{sec:geography}

\subsection{Cross-border ownership relations}

Among \NonrootEdges\ reported child--parent relations below level 0,
\CrossBorderEdges\ cross a standardized jurisdiction boundary. The pooled
share is \CrossBorderPct\pct, the equal-parent mean is
\CrossBorderEqualPct\pct, and the leave-one-parent-out range is
\CrossBorderLooLow--\CrossBorderLooHigh\pct. Recomputing with the raw supplied
country labels gives essentially the same result (Table~\ref{tab:robustness}).
This is one of the most definition-stable structural facts in the study.

The interpretation remains modest. A cross-border edge says only that the
reported child and immediate parent carry different country labels. It does
not reveal operating location, management, tax residence, legal control, or
the purpose of the relation.

\subsection{First-hop gateway amplification}

Table~\ref{tab:gateways} and Figure~\ref{fig:gateways} treat each observed
level-0 entity as a gateway and count normalized descendants below it. Dutch
gateways have a mean of \DutchGatewayMean\ descendants and median of
\DutchGatewayMedian. The parent-equal mean is \DutchGatewayEqualMean, and the
pooled leave-one-parent-out means range from \DutchGatewayLooLow\ to
\DutchGatewayLooHigh. Restricting to gateways without a single majority child
branch gives a mean of \DutchGatewayRobustMean, so the average is not generated
only by one immediate acquisition branch per gateway.

Mauritian gateways also show greater amplification than the United States
comparison, although their median is well below their mean. United States
gateways are numerous and span many parents, but typically have fewer
descendants. The comparison illustrates why resident prevalence and network
role must be measured separately.

\begin{table}[tbp]
\centering
\caption{Downstream amplification by selected first-hop jurisdiction}
\label{tab:gateways}
\scriptsize
\resizebox{\textwidth}{!}{\begin{tabular}{lrrrrrr}
\toprule
Gateway jurisdiction & Gateways & Parents & Mean & Median & Parent-equal mean & LOO mean range \\
\midrule
Netherlands & 13 & 12 & 34.8 & 12 & 37.1 & 21.2 to 39.7 \\
Mauritius & 17 & 11 & 17.9 & 5 & 19.7 & 12.4 to 19.5 \\
United States & 28 & 12 & 9.9 & 1.5 & 9.2 & 7.2 to 11.3 \\
Singapore & 18 & 12 & 7.5 & 3.5 & 7.4 & 6.0 to 8.5 \\
United Arab Emirates & 9 & 6 & 2.6 & 1 & 2.1 & 1.6 to 2.9 \\
\bottomrule
\end{tabular}
}
\begin{minipage}{0.97\textwidth}
\footnotesize\vspace{0.4em}
\textit{Notes:} One gateway is one observed level-0 normalized entity. Mean
and median weight gateways. The parent-equal mean first averages gateways
within parent. The LOO range repeats the gateway-weighted mean after omitting
each parent. The table shows selected jurisdictions; the full generated table
contains every jurisdiction with at least three gateways.
\end{minipage}
\end{table}

\begin{figure}[tbp]
\centering
\includegraphics[width=0.93\textwidth]{figure_04_gateway_amplification.pdf}
\caption{First-hop jurisdictions differ in downstream amplification}
\label{fig:gateways}
\begin{minipage}{0.91\textwidth}
\footnotesize\textit{Notes:} Each blue point is an observed level-0 gateway;
diamonds mark jurisdiction medians. The horizontal axis uses a symmetric-log
display to retain zero-descendant gateways. The plotted jurisdictions each
have at least five observed gateways.
\end{minipage}
\end{figure}

\subsection{Residence, intermediary prevalence, and path exposure}

There are \ObservedIntermediaries\ unique observed parent nodes in the logical
entity graph. Table~\ref{tab:roles} contrasts their distribution with resident
entity counts and target paths. The United States is the largest resident
location among the displayed jurisdictions. The Netherlands and Mauritius,
however, have substantially more target paths below them relative to their
resident counts. Figure~\ref{fig:roles} generalizes that comparison: points
above the diagonal have more descendant paths than resident entities, while
points below it are more terminal in this mapped topology.

The Netherlands is strictly upstream of \DutchUpstreamPct\pct\ of target
paths, compared with \DutchUpstreamEqualPct\pct\ after equal-parent weighting.
The pooled value is therefore partly a function of the sizes of the parent
networks using Dutch routes, but the role does not disappear when parents are
balanced.

The greedy union of the Netherlands, United States, and Mauritius is strictly
upstream of \TopThreeCount/\Targets\ target paths (\TopThreePct\pct). The
equal-parent union is only \TopThreeEqualPct\pct, and removing a large set of
venture-fund children gives \TopThreeNoFundsPct\pct. The gap is substantive:
the pooled cover is a descendant-weighted exposure statistic, not the share of
unique intermediaries and not a counterfactual measure of what would be
disconnected if three jurisdictions vanished.

\begin{table}[tbp]
\centering
\caption{Selected jurisdiction roles under distinct denominators}
\label{tab:roles}
\scriptsize
\resizebox{\textwidth}{!}{\begin{tabular}{lrrrrrr}
\toprule
Jurisdiction & Residents & Intermediaries & Paths below & Pooled (\%) & Parent-equal (\%) & Parents \\
\midrule
United States & 363 & 69.0 & 381 & 20.8 & 14.9 & 16.0 \\
Netherlands & 110 & 59.0 & 540 & 29.4 & 20.9 & 16.0 \\
Mauritius & 57.0 & 25.0 & 308 & 16.8 & 14.7 & 10.0 \\
United Kingdom & 99.0 & 39.0 & 219 & 11.9 & 10.2 & 16.0 \\
Singapore & 60.0 & 21.0 & 147 & 8.0 & 13.6 & 13.0 \\
\bottomrule
\end{tabular}
}
\begin{minipage}{0.97\textwidth}
\footnotesize\vspace{0.4em}
\textit{Notes:} Residents are parent-scoped normalized entities. Intermediaries
are unique observed parent nodes. Paths below count target paths on which the
jurisdiction appears strictly upstream, including a supplied country on a
named unresolved node. Pooled and parent-equal percentages use target paths;
``Parents'' is the number of parent buckets affected.
\end{minipage}
\end{table}

\begin{figure}[tbp]
\centering
\includegraphics[width=0.88\textwidth]{figure_05_jurisdiction_roles.pdf}
\caption{Resident prevalence and downstream network reach are distinct}
\label{fig:roles}
\begin{minipage}{0.88\textwidth}
\footnotesize\textit{Notes:} Both axes are logarithmic. Marker area increases
with the number of parent groups affected. Color is the jurisdiction's share
of unique observed intermediary nodes divided by its share of resident
entities, capped for display. India is excluded as the ultimate root.
\end{minipage}
\end{figure}

\subsection{Why a fixed ``centre'' category is not a core result}

A maintained list of frequently discussed financial or intermediary locations
is upstream of \DeclaredCentrePct\pct\ of target paths, and the equal-parent
estimate is \DeclaredCentreEqualPct\pct. But only
\DeclaredCentreNodePct\pct\ of unique observed intermediary nodes fall in the
same list. The path figure repeats a shared upstream node for every descendant,
and the list itself is imposed rather than observed. It also combines
locations that may host operating, financing, acquisition, or holding
entities. The canonical study therefore reports specific measured roles and
keeps the aggregate classification as a labeled sensitivity analysis.

\FloatBarrier

\section{Reported depth and observed graph distance}
\label{sec:depth}

Nearly half of targets are reported below the first foreign layer:
\DepthTwoPlusCount/\Targets\ (\DepthTwoPlusPct\pct) have reported level 2 or
higher. The apparently deep tail is less stable. Reported level 5 or higher
contains \DepthFivePlusCount\ targets (\DepthFivePlusPct\pct), while the
equal-parent mean is only \DepthFiveEqualPct\pct. Among complete paths, just
\GraphDepthFivePct\pct\ have reconstructed distance 5 or greater.

The measures disagree directly for \DepthMismatchCount/\CompletePaths
complete paths (\DepthMismatchPct\pct). Motherson contributes
\MothersonMismatchShare\pct\ of these mismatches. Figure~\ref{fig:depth}
shows that the maximum reported and reconstructed depths can diverge sharply
within a group. Reported level may preserve omitted layers; reconstructed
distance may be a lower bound. The concentration of disagreements in one
parent also means that a pooled depth correction cannot be assumed to apply
uniformly.

The sector result provides a falsification exercise. Of
\ManufacturingTargets\ targets labeled general manufacturing,
\DeepManufacturing\ are reported at level 5 or deeper. That looks like a broad
sector-depth association until parent composition is examined: Motherson and
Hindalco account for \DeepManufacturingTopTwo/\DeepManufacturing
(\DeepManufacturingTopTwoPct\pct) of the deep manufacturing targets. The
paper therefore rejects ``manufacturing subsidiaries are generally deeper''
as a sample-wide conclusion.

\begin{figure}[tbp]
\centering
\includegraphics[width=0.91\textwidth]{figure_06_reported_vs_graph_depth.pdf}
\caption{Reported depth and reconstructable graph distance diverge}
\label{fig:depth}
\begin{minipage}{0.91\textwidth}
\footnotesize\textit{Notes:} Each row is one parent bucket. Blue markers show
maximum supplied reported level. Orange markers show maximum reconstructed
distance among paths that reach the ultimate parent through observed entity
nodes. The connecting segment is descriptive; neither endpoint is designated
the legal truth.
\end{minipage}
\end{figure}

\FloatBarrier

\section{Ownership fields, financial coverage, and structural proxies}
\label{sec:financial}

\subsection{Recorded stakes and the ownership-depth gradient}

Stake information cannot support a complete control-weighted network. Every
observed level-0 target has missing mapping stake, and \ZeroStakeCount\ of
\NonrootEdges\ nonroot relations record exactly zero (\ZeroStakePct\pct).
Where a preferred financial row contains both mapping stake and an AOC-1
shareholding measure, each comparable zero-stake case is contradicted by a
positive AOC value (Appendix~\ref{app:ownership}). The zeros are therefore
treated as unknown, not as literal absence of ownership.

Restricting to positive-stake edges produces a superficially strong gradient:
the full-ownership rate is \OwnershipLevelOne\pct\ at reported level 1 and
\OwnershipLevelTwoPlus\pct\ at level 2 or deeper, a pooled difference of
\OwnershipPooledGap\ \pp. The equal-parent difference is only
\OwnershipEqualGap\ points, and excluding Reliance reduces the pooled gap to
\OwnershipNoRelianceGap\ points. The stable conclusion is a parent and
portfolio composition effect, not that ownership mechanically strengthens
with depth. Cumulative ownership along full paths is excluded because it would
condition on positive recorded stakes and treat unresolved zeros selectively.

\subsection{Financial attrition and extraction reuse}

Figure~\ref{fig:financial} and Table~\ref{tab:financial} document the financial
coverage funnel. A source is located for \SourceTargets/\Targets\ targets, a
balance sheet is parsed for \ParsedTargets, and only \ReadyTargets\ targets
are ready at least once. The preferred panel contains \ReadyRows\ ready
target-years; \SignPlausibleRows\ pass a positive-assets and nonnegative-
liabilities gate, and only \PlValidRows\ have a testable valid profit-and-loss
identity. Units are blank for \MissingUnitsReady/\ReadyRows\ ready rows.

Duplicate-evidence checks reinforce the limitation. There are
\RepeatedClusters\ repeated signatures defined jointly by fiscal year,
currency, units, assets, liabilities, and equity; they cover \RepeatedRows
parsed rows. In
\EvidenceReuseClusters\ clusters, the same source URL is attached to different
entity names and the core evidence text is byte-identical. A broad union of
repeated-signature and sign flags leaves \BroadUnflaggedRows\ ready rows, but
that remainder is a diagnostic, not a validated analysis sample.

Readiness itself varies greatly across parents, and \ParentsNoReady\ parents
have no ready target. The raw association between reported depth 3+ and
readiness is \FinancialRawDepthGap\ \pp; a parent fixed-effect coefficient is
\FinancialFeDepthGap\ points, with leave-one-parent-out estimates from
\FinancialLooLow\ to \FinancialLooHigh\ points. This is attenuation rather
than elimination, but selection, unit missingness, and extraction reuse make a
financial-performance interpretation indefensible.

\begin{figure}[tbp]
\centering
\includegraphics[width=\textwidth]{figure_07_financial_attrition.pdf}
\caption{Financial coverage is sparse and quality-selected}
\label{fig:financial}
\begin{minipage}{0.94\textwidth}
\footnotesize\textit{Notes:} Panel (a) counts structural targets at successive
coverage gates. Panel (b) counts preferred target-year rows. The broad screen
flags either a repeated numeric signature or an implausible sign combination;
its unflagged remainder has not been manually validated.
\end{minipage}
\end{figure}

\begin{table}[tbp]
\centering
\caption{Financial coverage and quality diagnostics}
\label{tab:financial}
\scriptsize
\resizebox{\textwidth}{!}{\begin{tabular}{lrrrl}
\toprule
Gate or diagnostic & Numerator & Denominator & Value & Unit \\
\midrule
Targets with any source & 847 & 1,834 & 46.2 & targets \\
Targets with parsed balance sheet & 444 & 1,834 & 24.2 & targets \\
Targets ready at least once & 265 & 1,834 & 14.4 & targets \\
Ready rows & 560 & 3,567 & 15.7 & preferred target-years \\
Sign-plausible ready rows & 521 & 560 & 93.0 & ready rows \\
P\&L-valid rows & 105 & 560 & 18.8 & ready rows \\
Ready rows with blank units & 516 & 560 & 92.1 & ready rows \\
Demonstrated same-URL/different-name reuse clusters & 33.0 & 51 & 64.7 & repeated-signature clusters \\
Ready rows unflagged by broad duplicate/sign screen & 460 & 560 & 82.1 & not a validated final sample \\
\bottomrule
\end{tabular}
}
\begin{minipage}{0.98\textwidth}
\footnotesize\vspace{0.4em}
\textit{Notes:} The value column is a percentage except for diagnostics whose
unit states otherwise. Repeated numeric signatures are fiscal-year/currency/
units/assets/liabilities/equity clusters. Demonstrated reuse additionally
requires the same URL, different target names, and byte-identical core evidence
text.
\end{minipage}
\end{table}

\subsection{A limited name-based structural proxy}

One supporting result survives a within-cell test. \HoldingNameCount\ targets
contain ``Holding(s),'' ``Holdco,'' or ``Investment(s)'' in the normalized
name. \HoldingNameHolderPct\pct\ of flagged targets hold at least one observed
child, compared with \OtherNameHolderPct\pct\ of other targets. Residualizing
both the outcome and name flag within parent-by-country cells gives a
\HoldingNameFeGap-point association. This suggests that names can be a useful
high-precision, low-recall structural proxy when edges are absent. It does not
verify a legal role, and external role labels and held-out-parent validation
would be required before using it as a general classifier.

\FloatBarrier

\section{Robustness, heterogeneity, and evidence classification}
\label{sec:robustness}

Figure~\ref{fig:robustness} and Table~\ref{tab:robustness} collect the main
robustness results on one scale. Three patterns are notably stable. The
cross-border edge share barely moves under equal-parent weighting, raw country
labels, or leave-one-parent-out checks. Modal-UIN concentration also remains
high under parent balance. Dutch gateway amplification remains large under
equal-parent weighting, omission of any one parent, and the no-majority-branch
restriction.

Other estimates change meaningfully. The three-country upstream union and
reported deep share fall under equal-parent weighting. The full-ownership
depth gap shrinks sharply under parent correction. Financial readiness retains
a small negative within-parent depth association, but its LOO range and severe
measurement attrition preclude substantive inference. The maintained-centre
exposure is much smaller when the unit is unique observed intermediary nodes
rather than descendant paths.

No formal probability sample underlies the \Parents\ parent buckets, so the
paper does not present sampling-error $p$-values as though they supported
population inference. Robustness here means sensitivity to measurement and
composition: unique entities versus repeated targets, pooled versus parent-
balanced weights, omission of any one parent, alternative graph definitions,
fund exclusions, and explicit missingness rules.

\begin{figure}[tbp]
\centering
\includegraphics[width=\textwidth]{figure_08_robustness_ladder.pdf}
\caption{Pooled estimates, parent balance, and definition sensitivity}
\label{fig:robustness}
\begin{minipage}{0.96\textwidth}
\footnotesize\textit{Notes:} Grey segments are leave-one-parent-out ranges
where applicable. Blue circles show pooled or primary estimates, orange
squares show equal-parent estimates, and red triangles show the labeled
alternative definition. Rows share a percent scale but not necessarily the
same observation unit; Table~\ref{tab:robustness} states each unit.
\end{minipage}
\end{figure}

\begin{landscape}
\begin{table}[p]
\centering
\caption{Canonical robustness matrix}
\label{tab:robustness}
\scriptsize
\resizebox{\linewidth}{!}{\begin{tabular}{llrrrrl}
\toprule
Estimand & Unit & Primary & Parent-equal/FE & LOO range & Alternative & Definition \\
\midrule
Cross-border nonroot edges & \% & 57.6 & 58.7 & 56.1 to 61.8 & 57.5 & Raw-label sensitivity \\
Modal UIN channel & \% of targets & 59.7 & 62.1 & 57.2 to 62.8 & 61.2 & Median parent \\
Largest observed subtree & median parent \% & 43.2 & -- & -- & 41.7 & Strict DAG dominator \\
Dutch gateway amplification & mean descendants & 34.8 & 37.1 & 21.2 to 39.7 & 37.8 & No majority child branch \\
Netherlands strictly upstream & \% of targets & 29.4 & 20.9 & 22.4 to 33.0 & -- & Descendant/path estimand \\
Netherlands + US + Mauritius upstream & \% of targets & 62.2 & 46.7 & 59.6 to 63.6 & 59.8 & Exclude venture-fund children \\
Reported depth 5+ & \% of targets & 18.6 & 7.3 & 8.6 to 20.9 & 7.4 & Reconstructed depth 5+, complete paths \\
Full-ownership depth gap & percentage points & 26.4 & 2.3 & -- & 6.1 & Exclude Reliance \\
Financial-ready deep gap & percentage points & -6.6 & -2.4 & -6.6 to -0.8 & -- & Parent FE in equal-parent column \\
\bottomrule
\end{tabular}
}
\begin{minipage}{0.97\linewidth}
\footnotesize\vspace{0.4em}
\textit{Notes:} ``Primary'' is pooled unless the row labels a median-parent
estimand. FE denotes a fixed-effect coefficient. LOO repeats the relevant
pooled or fixed-effect estimator after dropping each parent. Alternatives are
not always percentages of the same unit and should be read with the definition
column.
\end{minipage}
\end{table}
\end{landscape}

\section{Limitations and interpretation boundary}
\label{sec:limitations}

The study's limitations are part of its empirical result rather than a generic
disclaimer.

\begin{itemize}
  \item \textbf{Sample scope.} The parent buckets are supplied, not sampled
  from a documented universe. Results describe these groups and cannot
  establish prevalence across Indian multinational firms.
  \item \textbf{Entity identity.} Normalized names and countries are candidate
  identities. They cannot distinguish every alias, reorganization, joint
  exposure, or legally separate same-name company without registration
  identifiers.
  \item \textbf{Hierarchy completeness.} Named missing parents truncate paths,
  and reported levels disagree with observed graph distance. External
  ownership charts are required to adjudicate which representation is legally
  or historically correct.
  \item \textbf{UIN semantics.} The data establish mapping-channel
  concentration but do not validate the meaning of UIN substrings or their
  use as dates, office codes, or ownership types. Vintage claims are therefore
  excluded.
  \item \textbf{Jurisdiction labels.} A country label is not proof of operating
  location, management, tax residence, or entity function. The data cannot
  identify a tax, treaty, financing, acquisition, governance, or operating
  motive.
  \item \textbf{Ownership measurement.} Missing level-0 stakes and contradicted
  zero values prevent control weighting and cumulative-ownership inference.
  \item \textbf{Financial selection.} Sparse coverage, blank units, mixed
  currencies, invalid signs, limited P\&L identities, and demonstrated evidence
  reuse rule out population financial comparisons.
  \item \textbf{Time and current status.} The source combines reported
  structures and source years but does not validate a common current ownership
  date for every edge. The network is a reconstruction of the supplied report,
  not a claim about present legal control.
\end{itemize}

Accordingly, unusual topology is not evidence of wrongdoing. The paper does
not label any company, jurisdiction, route, or transaction fraudulent,
illegal, evasive, abusive, or improper. Legal, regulatory, institutional, and
tax interpretations require external evidence and are outside the identified
design.

\section{Conclusion}

The canonical reconstruction changes the object of study from a financial-
source panel into an explicit set of target occurrences, normalized entity
candidates, edges, and ownership paths. That correction is not cosmetic: it
reverses parent rankings and prevents long source histories from masquerading
as large corporate networks.

Within the supplied sample, parent architectures are genuinely heterogeneous.
Dominant mapping channels and large subtrees are common, but their magnitudes
depend on whether the unit is a UIN, observed subtree, named-node-inclusive
subtree, or strict graph dominator. The Netherlands and Mauritius occupy
intermediary and gateway roles that are not reducible to resident counts, and
the majority-cross-border-edge result is stable to parent reweighting. These
are defensible facts about the mapped topology.

The negative findings are equally important. Deep manufacturing is almost
entirely a two-parent pattern. The pooled ownership-depth gradient is largely
composition. Reported depth is not observed graph distance. The financial
layer cannot support comparative performance claims. A credible study of
corporate ownership networks therefore depends as much on refusing the wrong
denominator and interpretation as on finding a large statistic.

Future work can extend the evidence boundary only with new validated inputs:
official entity identifiers, dated ownership charts, documented UIN semantics,
verified stake fields, harmonized financial units, and a transparent sampling
frame. Until then, the contribution is a reproducible, falsifiable map of the
reported network structure in this dataset.

\clearpage
\appendix
\section{Exact estimands and counting rules}
\label{app:estimands}

\subsection{Units}

The pipeline distinguishes the following units throughout:

\begin{itemize}
  \item a source candidate row for source and extraction audits;
  \item a preferred target-year for financial readiness and accounting checks;
  \item a target occurrence for reported path exposure;
  \item a parent-scoped normalized entity for within-group architecture;
  \item a global normalized entity for cross-parent unique-company sensitivity;
  \item a nonroot occurrence edge for direct child--parent geography;
  \item a logical entity edge for graph topology;
  \item a complete observed target path for reconstructed distance;
  \item an observed level-0 gateway for descendant amplification; and
  \item an ultimate-parent bucket for equal-parent and omission checks.
\end{itemize}

No table silently substitutes raw rows for entities or target paths for unique
intermediaries. Percentages in generated tables always travel with a stated
numerator, denominator, or unit.

\subsection{Missing values and duplicates}

Blank country or name fields are never imputed through fuzzy similarity.
Named missing parents are kept as unresolved name--country nodes. Recorded
zero stakes are treated as unknown for control analyses. Structural fields are
required to be invariant within target identifier. Financial duplicates are
screened using exact signatures over fiscal year, currency, units, assets,
liabilities, and equity; demonstrated evidence reuse additionally requires a
shared source URL, different entity names, and identical core evidence text.

\subsection{Subtree and path denominators}

For parent $p$ with normalized entity set $V_p$ and candidate observed node
$v$, the observed-subtree share is

\[
  S_{pv}=\frac{|\operatorname{Desc}(v)\cap V_p|}{|V_p|-1}.
\]

The reported parent statistic is the maximum $S_{pv}$ over observed parent
nodes. The named-node-inclusive variant expands the candidate set to unresolved
named parents. The strict dominator variant removes $v$ and counts only entity
nodes that cease to be reachable from the ultimate root; alternate graph paths
therefore reduce its value.

For an upstream-jurisdiction result, each terminal target contributes at most
one indicator per jurisdiction, regardless of how often that jurisdiction
appears on its path. A shared ancestor still appears once for every terminal
target below it, which is why target-path exposure is reported separately from
unique intermediary-node prevalence.

\section{Depth audit}
\label{app:depth}

\begin{table}[H]
\centering
\caption{Reported depth, reconstructed distance, and parent concentration}
\scriptsize
\resizebox{\textwidth}{!}{\begin{tabular}{lrrrl}
\toprule
Depth diagnostic & Numerator & Denominator & Value (\%) & Unit \\
\midrule
Reported level 2+ & 871 & 1,834 & 47.5 & target occurrences \\
Reported level 5+ & 342 & 1,834 & 18.6 & target occurrences \\
Reported level 5+, equal-parent & -- & 28 & 7.3 & parent mean \\
Reconstructed distance 5+ & 125 & 1,687 & 7.4 & complete paths \\
Reported/reconstructed mismatch & 236 & 1,687 & 14.0 & complete paths \\
Motherson share of mismatches & 208 & 236 & 88.1 & mismatching complete paths \\
Manufacturing targets at reported level 5+ & 242 & 508 & 47.6 & manufacturing target occurrences \\
Motherson + Hindalco share of deep manufacturing & 240 & 242 & 99.2 & reported level-5+ manufacturing targets \\
\bottomrule
\end{tabular}
}
\begin{minipage}{0.98\textwidth}
\footnotesize\vspace{0.4em}
\textit{Notes:} Reported measures use all structural targets. Reconstructed
distance and the mismatch rate use only complete observed paths. The
equal-parent deep rate averages the parent-specific reported level-5+ shares.
\end{minipage}
\end{table}

\section{Ownership-field audit}
\label{app:ownership}

\begin{table}[H]
\centering
\caption{Stake coverage and ownership-depth diagnostics}
\scriptsize
\resizebox{\textwidth}{!}{\begin{tabular}{lrrrl}
\toprule
Ownership diagnostic & Numerator & Denominator & Value & Unit \\
\midrule
Level-0 stakes missing & 184 & 184 & 100 & level-0 targets \\
Nonroot recorded zero stakes & 406 & 1,650 & 24.6 & nonroot edges \\
Positive nonroot stakes & 1,244 & 1,650 & 75.4 & nonroot edges \\
Rows with mapping stake and AOC share & 45.0 & 3,567 & 1.3 & preferred target-years \\
Zero mapping stake but positive AOC share & 7.0 & 7 & 100 & comparable zero-stake rows \\
Stake/AOC difference above 1 pp & 8.0 & 45 & 17.8 & comparable rows \\
Full ownership at level 1 & 292 & 518 & 56.4 & positive-stake edges \\
Full ownership at level 2+ & 601 & 726 & 82.8 & positive-stake edges \\
Pooled level-2+ minus level-1 gap & -- & 1,244 & 26.4 & percentage points \\
Equal-parent level-2+ minus level-1 gap & -- & 27 & 2.3 & percentage points \\
Gap excluding Reliance & -- & 1,081 & 6.1 & percentage points \\
\bottomrule
\end{tabular}
}
\begin{minipage}{0.98\textwidth}
\footnotesize\vspace{0.4em}
\textit{Notes:} Full ownership is defined as a recorded positive stake of at
least 99.5 percent. Level-0 stakes are missing and nonroot zeros are not imputed.
The equal-parent gap averages parent-specific rates; groups without both depth
cells do not identify a paired contrast.
\end{minipage}
\end{table}

\section{Evidence disposition}
\label{app:evidence}

\begin{landscape}
\begin{table}[H]
\centering
\caption{Canonical classification of discovered and reviewed findings}
\scriptsize
\input{outputs/final/tables/paper_A3_evidence_classification.tex}
\begin{minipage}{0.94\linewidth}
\footnotesize\vspace{0.4em}
\textit{Notes:} ``Core'' requires agreement across the independent research
streams plus robustness to the relevant denominator, parent-weighting, and
hierarchy checks. Supporting results are defensible but not the paper's main
claim. Fragile, unresolved, and rejected claims are excluded from the core
narrative.
\end{minipage}
\end{table}
\end{landscape}

\section{Reproducibility}
\label{app:reproducibility}

From the repository root, the complete analysis and verification run with:

\begin{verbatim}
python -m src.analysis.run_pipeline --clean
python -m src.analysis.validate
pdflatex -interaction=nonstopmode -halt-on-error \
  -output-directory=paper paper/main.tex
pdflatex -interaction=nonstopmode -halt-on-error \
  -output-directory=paper paper/main.tex
\end{verbatim}

The analysis command checks immutable input hashes and exact central
invariants before writing any reported output. The validator rebuilds into a
temporary directory and byte-compares deterministic tables, PNG figures,
metrics, and generated research reports. Figure PDFs and the compiled paper
may contain run-time PDF metadata; their content is checked by compilation and
visual inspection rather than byte equality. Full definitions and the output
map appear in \texttt{REPLICATION.md}.


\end{document}
